\section{Sharp compatible bounds and extremizers}
\label{sec:compatible-bounds}

\subsection{The abelianization bound}

\begin{lemma}\label{lem:compatible-d}
A nonabelian compatible solution of class \(c\), with
\(d=\dim_{\F_4}(V/W)\), satisfies \(c\le d\).
\end{lemma}
\begin{proof}
Use a restricted square from Theorem~\ref{thm:compatibility}.
The identity
\(P(x)\beta(x,y)=P(y)r(x)\), with a lift of \(x\) understood,
implies
\begin{equation}\label{eq:double-endpoint}
 P(x)P(y)\beta(x,y)=0.
\end{equation}
In the universal module this kills the submodule \(K\) spanned by
\(t_it_jv_{ij}\), \(i<j\).
It has dimension \(\binom d2\) and lies in coefficient degree two.
Every multiplication by an endpoint kills it by a square, and
multiplication by a third index kills it by a triangle relation.

This relation space is invariant under \(\GL(U)\).
Indeed the universal expression
\(P(x)P(y)\beta(x,y)\) is additive in each input and scales
quadratically in each. For example its additive cross term in \(x\)
is
\[
 P(y)\bigl(P(x)\beta(z,y)+P(z)\beta(x,y)\bigr)
       =P(y)^2\beta(x,z)=0
\]
by Jacobi. Its basis values therefore span precisely \(K\).

Put \(\overline W_d=W_d/K\). The multigrading in
Proposition~\ref{prop:Hilbert} shows that its highest coefficient degree
is \(d-1\), for every \(d\ge2\), and that its top has the basis
\(T_j\) of \eqref{eq:terminal-basis}.
This top is \(\det(U)\otimes U\) also for \(d=2,3\): the same
diagonal and transvection calculations apply, with any additional
double-endpoint term now in \(K\). For \(d=2\) this is immediate
from \(T_1=t_1v_{12},T_2=t_2v_{12}\).

Every compatible derived module \(M\) is a quotient of \(\overline W_d\),
so \(c\le d+1\). Suppose \(c=d+1\). Choose a functional on \(M\)
nonzero on its terminal component. A change of quotient basis normalizes
its top pullback to \(\lambda(T_j)=1\), \(\lambda(T_k)=0\) for
\(k\ne j\).
Using \(P_j a_{ij}=P_i b_j\), we have, for \(i\ne j\),
\begin{equation}\label{eq:terminal-square-value}
 \im(T_j)=\Bigl(\prod_{h\ne j}P_h\Bigr)b_j.
\end{equation}
Write \(b_j=\sum c_{kl}a_{kl}\pmod{\m M}\), using intrinsic
generation. The product on the right kills \(\m M\).
If \(j\) is an endpoint of \(v_{kl}\), then
\((\prod_{h\ne j}t_h)v_{kl}\) is a top vector whose doubled
endpoint is different from \(j\). If \(j\) is not an endpoint,
the vector has both endpoints doubled and is in \(K\).
Thus the normalized top functional is zero on the right side of
\eqref{eq:terminal-square-value}, but one on the left.
This contradiction proves \(c\le d\).
\end{proof}

\subsection{Counting preimages of an operator cube}

\begin{lemma}\label{lem:compatible-w}
A compatible solution of class \(c\ge2\) satisfies
\[
 \dim_{\F_4}W\ge2^{c-1}-1.
\]
\end{lemma}
\begin{proof}
The case \(c=2\) says \(W\ne0\). Let \(c\ge3\).
By Lemma~\ref{lem:terminal-repeat}, choose a nonzero terminal value
\[
 P(x)P(z_1)\cdots P(z_{c-3})\beta(x,y).
\]
Restrictedness gives \(P(x)\beta(x,y)=P(y)b\), where
\(b=r(x)\in W\). Put \(m=c-2\),
\(u_1=y,u_2=z_1,\ldots,u_m=z_{c-3}\).
Then \(P(u_1)\cdots P(u_m)b\ne0\).
The \(2^m\) vectors
\[
 b_S=\Bigl(\prod_{i\in S}P(u_i)\Bigr)b
 \qquad(S\subseteq\{1,\ldots,m\})
\]
are independent by the complementary-operator argument used in
Theorem~\ref{thm:unrestricted-minimum}. Let \(C\) be their span.
Since \([x,r(x)]_0=0\), \(P(x)\) kills \(b\), hence all of \(C\).

For each nonempty \(S\), choose \(i\in S\) and set
\[
 v_S=\Bigl(\prod_{j\in S\setminus\{i\}}P(u_j)\Bigr)\beta(x,u_i).
\]
These vectors lie in the intrinsic derived ideal, and restrictedness gives
\(P(x)v_S=b_S\). Thus
\[
 \dim\ker(P(x)|_W)\ge2^m,\qquad
 \rank(P(x)|_W)\ge2^m-1.
\]
Rank-nullity yields \(\dim W\ge2^{m+1}-1=2^{c-1}-1\).
\end{proof}

\begin{theorem}\label{thm:compatible-minimum}
For every \(c\ge2\), the minimum binary dimension among compatible
solutions of polar class \(c\) is
\[
 \nu(c)=2^c+2c-2.
\]
\end{theorem}
\begin{proof}
Lemmas~\ref{lem:compatible-d} and \ref{lem:compatible-w} give the
lower bound by addition.
For attainment put \(s=c-1\),
\(R_s=\F_4[t_1,\ldots,t_s]/(t_i^2)\), and use
\begin{equation}\label{eq:semidirect-benchmark}
 f(x,z)=\left(\omega x,\omega z+\Bigl(\sum_i x_it_i\Bigr)z\right)
       \quad\text{on }\F_4^s\oplus R_s.
\end{equation}
The square of \(\sum x_it_i\) is zero, so direct substitution gives
\eqref{eq:functional}. The polar bracket is \(\F_4\)-bilinear and
has intrinsic derived ideal \(\m R_s\).
The scalar actions induced by \(f\) on that ideal and on the
abelianization are the displayed ones, so the map is compatible.
Its lower central ideals are
\(\gamma_k=\m^{k-1}R_s\) for \(k\ge2\), giving class \(s+1=c\).
The derived and quotient dimensions are \(2^s-1\) and \(s+1\).
The binary dimension is therefore \(2(s+2^s)=2^c+2c-2\).
\end{proof}

For \(c\ge5\), the extra four binary dimensions have an exact structural
meaning. The sharp bounds in the two subclasses are
\[
\begin{array}{c|cc}
 & d & w\\ \hline
\text{metabelian} & c-1 & 2^{c-1}-2\\
\text{compatible metabelian} & c & 2^{c-1}-1
\end{array}
\]
where each entry is the corresponding lower bound, attained
simultaneously. Compatibility therefore requires one additional
\(\F_4\)-dimension in the quotient and one in the derived ideal,
giving four additional binary dimensions. The restricted square
explains the strengthened bounds: it constrains terminal commutators
and supplies independent preimages for the nonempty faces of the
operator cube. The elementary semidirect construction attains this cost.

\subsection{Several extremizers in every class}

\begin{theorem}\label{thm:extremizers}
For every \(c\ge4\), at least \(c-2\) binary-linear conjugacy classes
attain \(\nu(c)\).
\end{theorem}
\begin{proof}
Set \(s=c-1\ge3\) and
\[
 R=\F_4[t_0,\ldots,t_{s-1}]/(t_i^2),\qquad W=\m R.
\]
Let \(U\) have basis \(x,u_1,\ldots,u_{s-1},v\).
Choose either \(h=0\), or \(h=t_1\cdots t_r\) with
\(1\le r\le s-2\), and put \(p=t_0h\).
Define the action on \(W\) by multiplication:
\[
 P_x=t_0,\qquad P_{u_i}=t_i,\qquad P_v=p.
\]
Define the alternating cocycle by
\[
 \beta(x,u_i)=t_i,\quad\beta(x,v)=t_0,\quad
 \beta(v,u_i)=(1+h)t_i,\quad\beta(u_i,u_j)=0.
\]
The only Jacobi triple types to check are \((x,u_i,u_j)\),
\((v,u_i,u_j)\), and \((x,v,u_i)\). Their sums are, respectively,
\[
 t_it_j+t_jt_i,\quad
 t_i(1+h)t_j+t_j(1+h)t_i,\quad
 t_0(1+h)t_i+pt_i+t_it_0,
\]
all zero. The cocycle values generate \(W\) under the action.

Take square values \(b_x=t_0,b_{u_i}=0,b_v=p\). They satisfy
\eqref{eq:compatibility-system}; the nontrivial relation involving \(p\)
uses \(p(1+h)=p\), since \(h^2=0\).
In scalar coordinates \(a=(X,Y_1,\ldots,Y_{s-1},V)\), the resulting
restricted square has the form
\begin{align}
 r(a,w)&=(0,P(a)w+q(a)),\label{eq:deformation-square}\\
 q(a)&=t_0(X^2+XV+hV^2)
       +\sum_i t_iY_i\bigl(X+(1+h)V\bigr).
       \label{eq:deformation-q}
\end{align}
All ring elements in these formulas are fixed coefficients. Consequently
\(f=\omega\id+r\) has binary degree at most two.
The square vanishes on \(W\), takes values in \(W\), and
\([a,r(a)]_0=0\). Expanding \(r(\omega a+r(a))\) proves
\(f^2=f+\id\). Its original polar is the displayed \(\F_4\)-bracket,
so it is compatible with the intrinsic scalar actions.

The lower central ideals are \(\gamma_k=\m^{k-1}R\) for \(k\ge2\),
since the action includes every \(t_i\). Each map has class \(c\)
and dimension \(\nu(c)\).

To distinguish them, use only intrinsic binary invariants. Let
\[
 \mathcal A=\im\bigl(L\longrightarrow\End_{\F_2}(W),\
                   a\longmapsto\ad a|_W\bigr),\qquad
 \mathcal A_k=\{D\in\mathcal A:D(W)\subseteq\gamma_{k+2}(L)\}.
\]
For \(h=0\), \(\dim_{\F_4}\mathcal A=s\) and
\(\mathcal A_k=0\) for \(k\ge2\).
For \(h=t_1\cdots t_r\), multiplication by \(p\) is independent of
the linear multiplications and is nonzero on \(W\), as \(r\le s-2\).
Thus \(\dim_{\F_4}\mathcal A=s+1\), and
\[
 \dim_{\F_4}\mathcal A_k=
 \begin{cases}1,&2\le k\le r+1,\\0,&k\ge r+2.\end{cases}
\]
Indeed any nonzero linear multiplication has a nonzero degree-two
output on \(W\), whereas \(pW\) has minimum nonzero degree \(r+2\).
The binary dimensions are twice these dimensions and are preserved by
every binary-linear conjugacy, independently of a choice of scalar lift.
They separate \(h=0\) and the \(s-2\) possible values of \(r\).
\end{proof}

The theorem gives an explicit collection of extremizers. It leaves open
the classification of all equality cases and their automorphism groups.
